\section{What a first failing index must satisfy}
\label{ext:smoothing}

The finite arguments use the ratio extremizer of
Corollary~\ref{ext:ratio-attainment}. A simple recurrence supplies the
same small excess at any first failure of the sharp bound. This lets us
use local exclusions directly, without first proving a separate sequence
of positive upper barriers.

\begin{lemma}[Gaussian derivative bounds]\label{ext:gaussian-constants}
The standard normal density satisfies
\[
 \|\phi''\|_\infty=\phi(0)<\frac25,
 \qquad \|\phi'''\|_\infty\le\frac2\pi,
 \qquad \frac1{4\pi}<\frac2{25},\qquad \CE<\frac{41}{100}.
\]
\end{lemma}
\begin{proof}
Since $\phi''(x)=(x^2-1)\phi(x)$, its absolute maximum is at zero:
the other candidates $x=\pm\sqrt3$ have value
$2e^{-3/2}\phi(0)<\phi(0)$. The absolutely convergent Gaussian
Fourier integral gives
\[
 |\phi'''(x)|\le\frac1{2\pi}
       \int_{\R}|t|^3e^{-t^2/2}\,dt=\frac2\pi.
\]
The bounds for $\phi(0)$ and $1/(4\pi)$ follow from $\pi>25/8$.
For the last comparison one may use
$\sqrt{10}<31623/10000$ and $\pi>314159/100000$, square the
positive sides, and compare rational numbers.
\end{proof}

\begin{theorem}[One-step smoothing]\label{ext:one-step}
For every integer $n\ge2$,
\begin{equation}\label{ext:smoothing-recurrence}
 C_n\le\sqrt{\frac n{n-1}}\,C_{n-1}
       +\frac{\sqrt n}{15(n-1)^{3/2}}
       +\frac{2\sqrt n}{25(n-1)^2}.
\end{equation}
\end{theorem}
\begin{proof}
Fix one standardized marginal $P$, write $v=n-1$, and put
$G_v(x)=\Phi(x/\sqrt v)$. Conditioning on the last summand yields
\[
 |\Pp(S_n\le H)-\E G_v(H-X)|\le C_{n-1}\rho/\sqrt v.
\]
Taylor expansion of the smooth Gaussian CDF, followed by expectation
under this same marginal, gives
\[
 |\E G_v(H-X)-G_v(H)-\tfrac12G_v''(H)|
 \le\frac{\phi(0)\rho}{6v^{3/2}}<\frac\rho{15v^{3/2}}.
\]
Here $\E X=0$ and $\E X^2=1$ cancel the first two Taylor terms.
For the change of variance, $\partial_vG_v=G_v''/2$ and
$\partial_v^2G_v=G_v^{(4)}/4$. Thus
\[
 G_{v+1}(H)-G_v(H)-\tfrac12G_v''(H)
 =\frac14\int_0^1(1-t)G_{v+t}^{(4)}(H)\,dt.
\]
Since $G_w^{(4)}(H)=w^{-2}\phi'''(H/\sqrt w)$, this remainder
has absolute value at most $1/(4\pi v^2)<2/(25v^2)$.
Combine the three estimates, multiply by $\sqrt n/\rho$, and use
$\rho\ge1$. The bound is uniform in $P$ and $H$ and treats both
signs of the discrepancy, so it passes to the two suprema defining $C_n$.
\end{proof}

\begin{corollary}[Entry from a sharp predecessor]
\label{ext:first-failure-entry}
If $n\ge40$ and $C_{n-1}\le\CE\le C_n$, then
\[
 0\le n(C_n-\CE)<\frac{1433}{4875}<\frac3{10},
 \qquad C_n<\frac{167}{400}<0.4226,
 \qquad n(C_{n-1}-C_n)_+=0.
\]
For $n\ge42$ the excess bound improves to $3003/10250$ and
the score bound to $73/175$.
For every standardized marginal $P$, the same-law predecessor loss
$n(\sqrt{n-1}\Delta_{n-1}(P)/\rho(P)-C_n)_+$ is zero.
At a first failing index $n\ge43$ above a sharp base, both deleted
constants $C_{n-1}$ and $C_{n-2}$ are below $0.4253$.
\end{corollary}
\begin{proof}
Put $s=\sqrt{n/(n-1)}$. The recurrence and
$n(s-1)=\{n/(n-1)\}/(s+1)$ give
\[
 n(C_n-\CE)\le\frac n{n-1}
 \left\{\frac{\CE}{s+1}+\frac{s}{15}
                  +\frac{2s}{25\sqrt{n-1}}\right\}.
\]
For $n\ge40$, $n/(n-1)\le40/39$, $\sqrt{n-1}>6$, and
$1<s<51/50$, since $39\cdot51^2-40\cdot50^2=1439>0$.
The right side is therefore less than
\[
 \frac{40}{39}\left(\frac{41}{200}
                  +\frac{17}{250}+\frac{17}{1250}\right)
 =\frac{1433}{4875}<\frac3{10}.
\]
This gives $C_n<0.41+0.3/40=167/400$.
For $n\ge42$, $n/(n-1)\le42/41$ and
$1<s<81/80$; the last comparison follows from
$41\cdot81^2-42\cdot80^2=201>0$.
Using $\CE<41/100$ and $s+1>2$, the right side is less than
\[
 \frac{42}{41}\left(\frac{41}{200}
                    +\frac{27}{400}+\frac{27}{2000}\right)
 =\frac{3003}{10250}<\frac3{10}.
\]
Consequently $C_n<0.41+0.3/42=73/175$.
Every law's predecessor ratio is bounded by $C_{n-1}\le\CE\le C_n$,
so both loss assertions follow. At a first failure the two deleted
indices are sharp and $\CE<0.41<0.4253$.
\end{proof}

\subsection{The loose ceilings needed after index twenty-four}
\label{ext:finite-loose-ceilings}
\begin{lemma}\label{ext:first-failure25}
If $n\ge25$ and $C_{n-1}\le\CE\le C_n$, then
\[
 n(C_n-\CE)<\frac{4177}{13824}<\frac{303}{1000},
 \qquad C_n<\frac{145873}{345600}<\frac{423}{1000}.
\]
For each marginal the same-law predecessor allowance is zero.
\end{lemma}
\begin{proof}
Set $s=\sqrt{n/(n-1)}$. Theorem~\ref{ext:one-step} gives
\[
 n(C_n-\CE)\le\frac{n}{n-1}
 \left(\frac{\CE}{s+1}+\frac{s}{15}
                       +\frac{2s}{25\sqrt{n-1}}\right).
\]
For $n\ge25$ one has $n/(n-1)\le25/24$,
$1<s<49/48$ and $\sqrt{n-1}>24/5$.
The latter strict comparisons follow from
$24\cdot49^2-25\cdot48^2=24>0$ and $24>576/25$.
Together with $\CE<41/100$, these imply
\[
 n(C_n-\CE)<\frac{25}{24}
  \left(\frac{41}{200}+\frac{49}{720}+\frac{49}{2880}\right)
 =\frac{4177}{13824}.
\]
Dividing by $n\ge25$ proves the score bound. The exact positive
comparison margins are
$303/1000-4177/13824=1459/1728000$ and
$423/1000-145873/345600=1579/1728000$.
Finally each marginal's predecessor score is at most $C_{n-1}\le\CE\le C_n$.
\end{proof}

In the original floor-$25$ initial and resource calculations, the own-score
ceiling is $434477/10^6$ and the retained predecessor allowance is $5/16$.
In the floor-$31$ calculations they are respectively $432313/10^6$ and
$31/100$. Both own-score ceilings exceed $423/1000$, by respectively
$11477/10^6$ and $9313/10^6$. Both allowances are nonnegative.
Thus Lemma~\ref{ext:first-failure25} supplies strictly stronger entry
information at a first failure. The native targets and numerical inputs
remain exactly as originally specified. The stronger information licenses
their weaker hypotheses; it does not change any certificate inequality.

Using a floor-$N$ resource calculation at an actual index $n>N$ additionally
requires its weighted invariant and its actual-index support hypotheses.
Lemma~\ref{ext:first-failure25} supplies only the scalar entry conditions;
it does not transfer support or the bound $D\le d$ from one index to another.
The transfer used here is the proved $nD\le Nd$ from the common finite
resource calculus, together with separately verified actual-index support.
